\documentclass[a4paper,12pt]{article}
\usepackage[UTF8]{ctex}
\usepackage{graphicx}
\usepackage{geometry}
\usepackage{fancyhdr}
\usepackage{titlesec}
\usepackage{enumitem}
\usepackage{booktabs}
\usepackage{longtable}
\usepackage{array}
\usepackage{multirow}
\usepackage{amsmath}
\usepackage{hyperref}

% 页面设置
\geometry{top=3.0cm,bottom=3.0cm,left=2.5cm,right=2.5cm,headheight=2.0cm,footskip=2.0cm}

% 页眉设置
\pagestyle{fancy}
\fancyhf{}
\fancyhead[C]{西南科技大学设计概要}
\renewcommand{\headrulewidth}{0.8pt}

% 标题格式设置
\titleformat{\section}{\centering\zihao{-2}\heiti}{第\,\thesection\,章}{1em}{}
\titleformat{\subsection}{\zihao{-3}\heiti}{\thesubsection}{1em}{}
\titleformat{\subsubsection}{\zihao{4}\heiti}{\thesubsubsection}{1em}{}

\begin{document}

% 封面
\begin{titlepage}
\begin{center}
\vspace*{2cm}
{\zihao{1}\heiti 西南科技大学}\\[1cm]
{\zihao{1}\heiti 设计概要}\\[2cm]
{\zihao{2}\heiti 基于Pygame的2D RPG游戏设计与实现}\\[0.5cm]
{\zihao{3}\heiti ——以《西游记：祸起观音院》为例}\\[3cm]
\begin{tabular}{rl}
{\zihao{4}\songti 班\quad 级：} & {\zihao{4}\songti 软件2401} \\
{\zihao{4}\songti 学\quad 号：} & {\zihao{4}\songti 5120243335} \\
{\zihao{4}\songti 姓\quad 名：} & {\zihao{4}\songti 袁杰强} \\
\end{tabular}\\[2cm]
{\zihao{4}\songti \today}
\end{center}
\end{titlepage}

% 正文
\newpage
\section{绪论}

\subsection{课题背景与意义}
随着计算机技术和游戏产业的快速发展，电子游戏已经成为人们日常娱乐的重要组成部分。2D角色扮演游戏（RPG）因其独特的叙事方式和游戏体验，一直拥有广泛的玩家群体。中国古典文学作品如《西游记》蕴含着丰富的文化内涵和故事素材，将其与现代游戏技术相结合，不仅能够传承和弘扬传统文化，还能为玩家提供具有中国特色的游戏体验。

Pygame作为Python语言的游戏开发库，具有简单易用、跨平台、开源免费等优点，非常适合用于2D游戏的开发和教学研究。基于Pygame开发《西游记：祸起观音院》游戏，对于学习游戏开发技术、探索传统文化与现代技术的结合具有重要的实践意义。

\subsection{国内外研究现状}
在游戏开发领域，2D RPG游戏的开发技术已经相当成熟。国外以Unity、Unreal Engine等商业引擎为主导，国内则有Cocos2d-x等开源框架。Python语言凭借其简洁的语法和丰富的库支持，在游戏开发教育领域得到广泛应用。Pygame作为Python最流行的游戏开发库之一，提供了图形渲染、音频处理、事件管理等完整功能，被广泛用于游戏开发教学和独立游戏制作。

在传统文化与游戏结合方面，国内外都有成功案例。《黑神话：悟空》等作品证明了中国传统文化题材在游戏市场中的巨大潜力。然而，针对教学和研究目的的轻量级传统文化游戏项目相对较少，本项目填补了这一空白。

\subsection{本文研究的主要内容}
本文主要研究内容包括：
（1）分析2D RPG游戏的功能需求，设计合理的系统架构；
（2）实现基于Pygame的游戏引擎核心功能，包括场景管理、角色控制、碰撞检测等；
（3）设计并实现回合制战斗系统，支持多种攻击方式和技能组合；
（4）实现完整的NPC交互系统，包括对话、任务触发、属性增益等功能；
（5）设计测试用例，验证系统的功能完整性和稳定性。

\section{系统需求分析}

\subsection{功能需求}
根据游戏设计目标，系统需要实现以下功能：
\begin{enumerate}[label=(\arabic*)]
\item 游戏流程控制：包括游戏启动、场景切换、游戏结束等完整流程控制；
\item 角色控制：实现玩家角色的移动、加速、方向控制等功能；
\item 场景渲染：支持多场景（开始界面、村庄、寺庙）的渲染和切换；
\item NPC交互：实现与不同类型NPC的对话和交互功能；
\item 战斗系统：实现回合制战斗，支持普通攻击、技能攻击和大招；
\item UI界面：提供对话框、确认框、避战选择框等用户界面组件；
\item 音频支持：实现背景音乐和音效的播放；
\item 地图系统：支持Tiled地图编辑器生成的TMX格式地图。
\end{enumerate}

\subsection{非功能需求}
\begin{enumerate}[label=(\arabic*)]
\item 性能要求：游戏运行帧率不低于30FPS，响应延迟不超过100ms；
\item 兼容性要求：支持Windows、Linux、macOS等主流操作系统；
\item 可维护性：代码结构清晰，模块化设计，便于后续维护和扩展；
\item 可扩展性：支持新场景、新角色、新技能的添加。
\end{enumerate}

\section{系统设计}

\subsection{系统架构设计}
本系统采用分层架构设计，将系统划分为五个层次：
\begin{enumerate}[label=(\arabic*)]
\item 配置层（config）：负责全局常量、游戏参数、资源路径等配置管理；
\item 核心层（core）：提供游戏引擎核心功能，包括主控类、动画系统、地图系统；
\item 场景层（scenes）：实现各个游戏场景的逻辑和渲染；
\item 角色层（characters）：管理玩家角色和NPC的行为和属性；
\item 界面层（ui）：提供各种用户界面组件。
\end{enumerate}

\subsection{模块设计}

\subsubsection{配置模块（config）}
配置模块负责管理游戏的所有全局配置参数，包括：
\begin{itemize}
\item 窗口设置：屏幕尺寸（800×600）、帧率（30FPS）；
\item 玩家属性：移动速度、基础生命值（100）、基础攻击力（25）；
\item 战斗参数：普通攻击伤害（25）、技能伤害（40）、大招伤害（70）；
\item 怪物属性：普通怪物生命值（100）、BOSS生命值（250）；
\item 资源路径：图片、音频、地图、字体等资源文件路径。
\end{itemize}

\subsubsection{核心模块（core）}
核心模块包含三个主要组件：

（1）游戏主控类（Game）：负责游戏初始化、主循环控制、场景调度和玩家状态管理。采用状态机模式管理游戏流程，定义了四个状态常量：STATE\_START（开始）、STATE\_VILLAGE（村庄）、STATE\_TEMPLE（寺庙）、STATE\_END（结束）。

（2）动画系统（Animation）：提供帧动画管理功能，支持多角色动画资源加载。包括孙悟空4方向行走帧、牛妖多状态帧、土地公帧、战斗动画帧等加载函数。

（3）地图系统（TiledMap）：实现Tiled地图编辑器生成的TMX格式地图解析和渲染。支持对象组提取、道路多边形碰撞、障碍物矩形碰撞等功能。

\subsubsection{场景模块（scenes）}
场景模块实现四个主要场景：

（1）开始界面（StartScene）：显示游戏标题、开始按钮和退出按钮，提供游戏入口。

（2）村庄场景（VillageScene）：实现村庄地图的渲染和探索，包含玩家移动、NPC交互、场景过渡等功能。支持与土地公、长老、农民等NPC的对话和交互。

（3）寺庙场景（TempleScene）：实现寺庙地图的渲染和探索，包含怪物遭遇、战斗触发、避战选择等功能。支持与牛妖的战斗和避战交互。

（4）战斗场景（FightScene）：实现完整的回合制战斗系统，包含战斗状态机、攻击动画、伤害计算、胜负判定等功能。

\subsubsection{角色模块（characters）}
角色模块采用继承层次设计：

（1）玩家类（Player）：继承自pygame.sprite.Sprite，实现4方向移动、加速控制、动画播放、碰撞检测等功能。

（2）NPC基类（NPC）：提供NPC的通用功能，包括徘徊运动、面向目标、动画更新等。

（3）土地公类（God）：继承自NPC，实现限域徘徊、救援对话、冷却控制等功能。

（4）农民类（Farmer）：继承自NPC，实现轮流对话、冷却控制等功能。

（5）长老类（Elder）：继承自NPC，实现固定位置、增益对话、属性提升等功能。

（6）牛妖类（Cattle）：继承自NPC，实现多动画状态、HP/ATK属性、boss标记等功能。

\subsubsection{界面模块（ui）}
界面模块提供四种UI组件：

（1）对话框（Dialog）：显示文本对话，支持自动换行和图片背景。

（2）确认对话框（ConfirmDialog）：显示是/否选择，用于场景切换确认。

（3）避战选择框（AvoidBattleDialog）：显示战斗/避战选择，用于怪物遭遇交互。

（4）场景过渡（FadeScene）：实现场景渐入渐出效果，提供平滑的场景切换。

\section{接口设计}

\subsection{内部接口}
系统内部接口主要体现在场景间的数据传递：

（1）Game类维护全局玩家状态：player\_hp、player\_max\_hp、player\_atk、ultimate\_unlocked、rescued。

（2）场景创建时接收当前玩家状态：VillageScene(player\_hp, player\_max\_hp, ultimate\_unlocked, rescued)。

（3）场景结束时回写更新后的状态：Game类通过\_sync\_player\_state\_from\_village()方法同步状态。

\subsection{外部接口}
系统外部接口包括：

（1）文件系统接口：读取TMX地图文件、图片资源、音频资源、字体文件。

（2）Pygame接口：使用Pygame提供的图形渲染、事件处理、音频播放等功能。

（3）用户输入接口：处理键盘输入（方向键、数字键、功能键）和鼠标输入。

\section{数据结构设计}

\subsection{游戏状态数据}
\begin{table}[h]
\centering
\caption{游戏状态数据结构}
\begin{tabular}{lll}
\toprule
数据项 & 类型 & 说明 \\
\midrule
player\_hp & int & 玩家当前生命值 \\
player\_max\_hp & int & 玩家最大生命值 \\
player\_atk & int & 玩家攻击力 \\
ultimate\_unlocked & bool & 大招是否解锁 \\
rescued & bool & 是否被救援过 \\
current\_state & int & 当前游戏状态 \\
\bottomrule
\end{tabular}
\end{table}

\subsection{角色数据}
\begin{table}[h]
\centering
\caption{角色数据结构}
\begin{tabular}{lll}
\toprule
数据项 & 类型 & 说明 \\
\midrule
map\_x, map\_y & float & 地图坐标 \\
dx, dy & int & 速度分量 \\
direction & int & 当前方向 \\
hp, max\_hp & int & 生命值 \\
atk & int & 攻击力 \\
speed & int & 移动速度 \\
is\_boss & bool & 是否为BOSS \\
\bottomrule
\end{tabular}
\end{table}

\subsection{战斗状态数据}
\begin{table}[h]
\centering
\caption{战斗状态数据结构}
\begin{tabular}{lll}
\toprule
状态 & 值 & 说明 \\
\midrule
STATE\_INTRO & -1 & 开场喊话 \\
STATE\_READY & 0 & 等待玩家输入 \\
STATE\_PLAYER\_ATK & 1 & 玩家攻击动画 \\
STATE\_MONSTER\_HURT & 2 & 怪物受伤 \\
STATE\_MONSTER\_ATK & 3 & 怪物攻击动画 \\
STATE\_PLAYER\_HURT & 4 & 玩家受伤 \\
STATE\_WIN & 5 & 胜利 \\
STATE\_LOSE & 6 & 失败 \\
\bottomrule
\end{tabular}
\end{table}

\section{用例测试设计}

\subsection{测试环境}
\begin{itemize}
\item 操作系统：Windows 10/11
\item Python版本：3.6+
\item Pygame版本：2.0+
\item 测试工具：手动测试
\end{itemize}

\subsection{测试用例}

\subsubsection{游戏启动测试}
\begin{longtable}{|p{1.5cm}|p{3cm}|p{4cm}|p{4cm}|p{2cm}|}
\hline
编号 & 测试项 & 测试步骤 & 预期结果 & 实际结果 \\
\hline
TC001 & 游戏启动 & 双击main.py运行游戏 & 游戏窗口正常显示，显示开始界面 & 通过 \\
\hline
TC002 & 开始按钮点击 & 点击"开始游戏"按钮 & 进入村庄场景 & 通过 \\
\hline
TC003 & 退出按钮点击 & 点击"退出游戏"按钮 & 游戏正常退出 & 通过 \\
\hline
TC004 & 回车键开始 & 按回车键 & 进入村庄场景 & 通过 \\
\hline
TC005 & ESC键退出 & 按ESC键 & 游戏正常退出 & 通过 \\
\hline
\end{longtable}

\subsubsection{角色移动测试}
\begin{longtable}{|p{1.5cm}|p{3cm}|p{4cm}|p{4cm}|p{2cm}|}
\hline
编号 & 测试项 & 测试步骤 & 预期结果 & 实际结果 \\
\hline
TC006 & 方向键移动 & 按方向键控制角色移动 & 角色向对应方向移动 & 通过 \\
\hline
TC007 & 加速移动 & 按住Shift键移动 & 角色移动速度加倍 & 通过 \\
\hline
TC008 & 边界检测 & 移动到地图边界 & 角色无法超出地图边界 & 通过 \\
\hline
TC009 & 障碍物碰撞 & 移动到障碍物位置 & 角色无法穿过障碍物 & 通过 \\
\hline
TC010 & 动画播放 & 移动时观察角色动画 & 角色播放对应方向行走动画 & 通过 \\
\hline
\end{longtable}

\subsubsection{NPC交互测试}
\begin{longtable}{|p{1.5cm}|p{3cm}|p{4cm}|p{4cm}|p{2cm}|}
\hline
编号 & 测试项 & 测试步骤 & 预期结果 & 实际结果 \\
\hline
TC011 & 土地公对话 & 靠近土地公触发对话 & 显示对话框，展示对话内容 & 通过 \\
\hline
TC012 & 前往寺庙确认 & 土地公对话后确认前往 & 弹出确认框，选择"是"进入寺庙 & 通过 \\
\hline
TC013 & 长老1增益 & 与长老1对话 & HP上限+50，HP+50 & 通过 \\
\hline
TC014 & 长老2增益 & 与长老2对话 & 解锁大招"七十二变" & 通过 \\
\hline
TC015 & 长老3增益 & 与长老3对话 & HP完全恢复 & 通过 \\
\hline
TC016 & 农民对话 & 与农民对话 & 显示叙事对话内容 & 通过 \\
\hline
TC017 & 对话冷却 & 连续与同一NPC对话 & 存在冷却时间，无法连续对话 & 通过 \\
\hline
\end{longtable}

\subsubsection{战斗系统测试}
\begin{longtable}{|p{1.5cm}|p{3cm}|p{4cm}|p{4cm}|p{2cm}|}
\hline
编号 & 测试项 & 测试步骤 & 预期结果 & 实际结果 \\
\hline
TC018 & 怪物遭遇 & 靠近牛妖触发战斗 & 弹出避战选择框 & 通过 \\
\hline
TC019 & 选择战斗 & 在避战选择框选择"战斗" & 进入回合制战斗场景 & 通过 \\
\hline
TC020 & 选择避战 & 在避战选择框选择"避战" & 扣除20HP，返回场景 & 通过 \\
\hline
TC021 & 普通攻击 & 战斗中按数字键1 & 对怪物造成25点伤害 & 通过 \\
\hline
TC022 & 技能攻击 & 战斗中按数字键2 & 对怪物造成40点伤害 & 通过 \\
\hline
TC023 & 大招攻击 & 战斗中按数字键3（已解锁） & 对怪物造成70点伤害 & 通过 \\
\hline
TC024 & 怪物攻击 & 玩家攻击后怪物反击 & 对玩家造成怪物攻击力伤害 & 通过 \\
\hline
TC025 & 战斗胜利 & 怪物HP降为0 & 显示胜利信息，返回场景 & 通过 \\
\hline
TC026 & 战斗失败 & 玩家HP降为0 & 标记rescued，土地公救援 & 通过 \\
\hline
\end{longtable}

\subsubsection{场景切换测试}
\begin{longtable}{|p{1.5cm}|p{3cm}|p{4cm}|p{4cm}|p{2cm}|}
\hline
编号 & 测试项 & 测试步骤 & 预期结果 & 实际结果 \\
\hline
TC027 & 村庄到寺庙 & 确认前往寺庙 & 场景渐出后进入寺庙场景 & 通过 \\
\hline
TC028 & 寺庙返回村庄 & 在寺庙按ESC键 & 确认后返回村庄场景 & 通过 \\
\hline
TC029 & 战败返回 & 战斗失败后 & 土地公救援对话后返回村庄 & 通过 \\
\hline
TC030 & 通关显示 & 击败全部怪物 & 显示通关画面和恭喜信息 & 通过 \\
\hline
\end{longtable}

\subsubsection{UI组件测试}
\begin{longtable}{|p{1.5cm}|p{3cm}|p{4cm}|p{4cm}|p{2cm}|}
\hline
编号 & 测试项 & 测试步骤 & 预期结果 & 实际结果 \\
\hline
TC031 & 对话框显示 & 触发对话 & 对话框正常显示，文字清晰 & 通过 \\
\hline
TC032 & 对话框关闭 & 按空格/回车/ESC & 对话框正常关闭 & 通过 \\
\hline
TC033 & 确认框显示 & 触发确认操作 & 确认框正常显示，按钮可点击 & 通过 \\
\hline
TC034 & 确认框选择 & 点击"是"或"否" & 根据选择执行相应操作 & 通过 \\
\hline
TC035 & 场景过渡 & 场景切换时 & 渐入渐出效果正常显示 & 通过 \\
\hline
\end{longtable}

\subsubsection{音频系统测试}
\begin{longtable}{|p{1.5cm}|p{3cm}|p{4cm}|p{4cm}|p{2cm}|}
\hline
编号 & 测试项 & 测试步骤 & 预期结果 & 实际结果 \\
\hline
TC036 & 背景音乐 & 进入各场景 & 对应场景背景音乐正常播放 & 通过 \\
\hline
TC037 & 战斗音效 & 战斗中攻击 & 攻击音效正常播放 & 通过 \\
\hline
TC038 & 音乐切换 & 场景切换时 & 背景音乐平滑切换 & 通过 \\
\hline
\end{longtable}

\section{结论}

本文设计并实现了一款基于Pygame的2D RPG游戏《西游记：祸起观音院》。系统采用分层架构设计，包含配置层、核心层、场景层、角色层和界面层五个层次，具有良好的模块化和可扩展性。

通过详细的测试用例验证，系统各项功能运行正常，包括游戏启动、角色移动、NPC交互、回合制战斗、场景切换、UI组件和音频系统等。测试结果表明，系统满足设计要求，能够提供流畅的游戏体验。

本项目对于学习Pygame游戏开发、理解面向对象设计思想、探索传统文化与现代技术结合具有一定的参考价值。后续可以进一步扩展游戏内容，如增加更多场景、角色、技能和任务系统。

\end{document}
